\chapter{Financial Stress}

\section*{Definition and Overview}
Financial stress is the worry and strain that comes from difficulty paying bills, managing debt, or feeling financially insecure. It is one of the most common and powerful sources of stress in modern life, and it affects not only money but also sleep, relationships, mental health, and physical health. Financial stress can be triggered by job loss, illness, rising costs, relationship breakdown, or debt, and it often becomes a vicious cycle: stress impairs decision-making, which worsens finances, which increases stress. The good news is that practical steps, support services, and proven strategies can reduce financial stress, and help is available for free. This chapter explains the effects of financial stress and how to manage it.

\section*{Epidemiology}
Financial stress affects a large proportion of the population. Surveys in Australia consistently find that money worries are among the top sources of stress, affecting around a third to half of adults at any given time, and rates rise sharply in economic downturns and with rising cost of living. Financial stress is more common in single-parent families, people with low or unstable income, renters, and people with chronic illness or disability. Money stress is strongly linked to anxiety, depression, relationship conflict, and even physical illness, and it is a recognised risk factor for suicide. Help-seeking for financial problems is often delayed by shame, which worsens outcomes.

\section*{Aetiology and Pathophysiology}
\begin{itemize}
    \item \textbf{Triggers:} job loss, reduced income, unexpected expenses, medical bills, relationship breakdown, and rising costs of housing, food, and energy
    \item \textbf{Debt:} credit cards, personal loans, and mortgage or rent arrears create ongoing pressure, especially with high interest
    \item \textbf{Income inequality:} low and insecure income makes it hard to build savings or absorb shocks
    \item \textbf{The stress response:} financial threat activates the same fight-or-flight response as physical danger, releasing cortisol and adrenaline
    \item \textbf{Health effects:} chronic stress disturbs sleep, raises blood pressure, weakens immunity, and increases the risk of anxiety, depression, and heart disease
    \item \textbf{Behavioural effects:} stress impairs decision-making, leading to avoidance, impulsive spending, or "present bias" that worsens finances
    \item \textbf{Social effects:} shame and stigma cause withdrawal, relationship conflict, and reduced help-seeking
    \item \textbf{The cycle:} financial stress harms health; poor health and chronic illness then reduce earning capacity, deepening the stress
\end{itemize}

\section*{Clinical Features}
\begin{itemize}
    \item Persistent worry about money, bills, and the future
    \item Difficulty sleeping, waking with racing thoughts about finances
    \item Irritability, low mood, anxiety, and reduced concentration
    \item Physical symptoms: headaches, muscle tension, fatigue, and digestive problems
    \item Relationship conflict about money
    \item Avoidance: not opening bills, ignoring debt letters, and delaying financial decisions
    \item Behavioural changes: cutting back on essentials, borrowing, or using credit to cover basic needs
    \item Effects on work: reduced performance, presenteeism, and sickness absence
\end{itemize}

\section*{Differential Diagnosis}
Financial stress often coexists with mental health conditions that also need attention.
\begin{itemize}
    \item Generalised anxiety disorder and depression, which may be worsened by or independent of money worries
    \item Adjustment disorder after a financial shock such as job loss
    \item Problem gambling or substance use, which can both cause and result from financial stress
    \item Other causes of insomnia and physical symptoms, which should be assessed
    \item The distinction matters because treating the mental health condition helps the person act on the financial plan
\end{itemize}

\section*{When to Seek Medical Care}
Anyone whose financial stress is causing persistent sleep problems, low mood, anxiety, or physical symptoms should see a doctor, who can also connect them with support services. Free financial counselling is available to everyone regardless of income.

\textbf{Call 000 (or local emergency services) for:}
\begin{itemize}
    \item Suicidal thoughts or behaviour, which are more common in people under severe financial stress
    \item Chest pain, severe breathlessness, or collapse
    \item Any crisis in which someone cannot keep themselves safe
\end{itemize}

\textbf{Support lines:} Lifeline 13 11 14 and the Suicide Call Back Service 1300 659 467 are available around the clock.

\section*{Diagnosis and Investigations}
\begin{itemize}
    \item \textbf{History:} the doctor explores the stress, its triggers, and its effects on sleep, mood, and functioning
    \item \textbf{Mental health screening:} questionnaires for anxiety and depression
    \item \textbf{Physical review:} blood pressure and examination where stress-related symptoms are present
    \item \textbf{Financial assessment:} a free financial counsellor reviews income, expenses, debts, and options in a confidential setting
    \item \textbf{Social work referral:} for help with housing, utilities, and other practical matters
\end{itemize}

\section*{Management}
\begin{itemize}
    \item \textbf{Free financial counselling:} National Debt Helpline services help people negotiate with creditors, set up payment plans, and manage debt without judgement
    \item \textbf{Make a budget:} writing down income and expenses makes the situation concrete and controllable, and small wins build confidence
    \item \textbf{Prioritise essentials:} sort expenses into needs (housing, food, energy) and wants, and address arrears before they escalate
    \item \textbf{Seek entitlements:} Centrelink payments, concessions, energy rebates, and other assistance are often underclaimed
    \item \textbf{Address debt:} negotiate with creditors, consolidate where helpful, and consider formal options through a financial counsellor
    \item \textbf{Mental health care:} psychological therapy and, where needed, medication for anxiety and depression
    \item \textbf{Stress management:} sleep routines, exercise, relaxation, and limits on how much time is spent worrying about money
    \item \textbf{Employment support:} job-seeking services and retraining pathways after job loss
    \item \textbf{Support networks:} talking to trusted people and community services reduces shame and isolation
\end{itemize}

\section*{Complications}
\begin{itemize}
    \item Anxiety, depression, and other mental illness
    \item Insomnia and physical illness, including high blood pressure and heart disease
    \item Relationship breakdown and family conflict
    \item Problem gambling, substance use, and impulsive spending as coping attempts
    \item Homelessness, eviction, and disconnection of utilities when arrears are not addressed
    \item Suicide, the most serious risk
\end{itemize}

\section*{Prognosis and Outcomes}
Financial stress can be overcome. Free financial counselling helps most people who use it to reduce debt, stabilise their situation, and build a workable plan, and the mental health effects respond to the same treatments that help other forms of stress. People who combine practical financial action with self-care and support recover better and faster. The key shift is from avoidance and shame to action: contacting a free counsellor, making a plan, and taking one small step at a time changes both the finances and the stress.

\section*{Prevention and Self-Care}
\begin{itemize}
    \item Open bills and letters promptly, and seek help before debts escalate
    \item Call the National Debt Helpline (1800 007 007) for free, confidential advice
    \item Make a simple budget and review it monthly
    \item Build a small emergency buffer where possible, even if it is modest
    \item Protect sleep, exercise, and time with supportive people
    \item Limit alcohol and avoid gambling as ways of coping
    \item Talk to a doctor about the mental health effects of money stress
    \item Accept that financial difficulty happens to many people and that seeking help is a strength
\end{itemize}

\section*{Sources and Further Reading}
The following reputable sources provide further information for healthcare students and patients seeking reliable, current advice about financial stress.
\begin{itemize}
    \item National Debt Helpline. \textit{Free financial counselling.} https://ndh.org.au/
    \item Services Australia. \textit{Payments and concessions.} https://www.servicesaustralia.gov.au/
    \item MoneySmart. \textit{Managing money and debt.} https://moneysmart.gov.au/
    \item Beyond Blue. \textit{Financial stress and mental health.} https://www.beyondblue.org.au/
    \item Lifeline. \textit{Crisis support.} https://www.lifeline.org.au/
    \item World Health Organization. \textit{Mental health and social determinants.} https://www.who.int/
\end{itemize}
